% Part: lambda-calculus
% Chapter: lambda-definability
% Section: introduction

\documentclass[../../../include/open-logic-section]{subfiles}

\begin{document}

\olfileid[id]{lam}{ldf}{int}
\olsection{Pendahuluan}

Sekilas, kalkulus lambda hanyalah kalkulus yang sangat abstrak bagi
ekspresi-ekspresi yang merepresentasikan fungsi serta penerapan fungsi-fungsi
itu pada ekspresi lain. Tidak ada sesuatu pun dalam sintaksis kalkulus lambda
yang menunjukkan bahwa fungsi-fungsi tersebut bekerja pada jenis objek
tertentu; khususnya, sintaksisnya sama sekali tidak menyebut bilangan asli.
Operasi-operasi dasarnya---aplikasi dan abstraksi lambda---berlaku pada setiap
fungsi, bukan hanya pada fungsi atas bilangan asli.

Meskipun demikian, dengan sedikit kecerdikan, fungsi aritmetika, yakni fungsi
atas bilangan asli, dapat didefinisikan dalam kalkulus lambda. Untuk
melakukannya, bagi setiap bilangan asli~$n \in \Nat$, kita definisikan suatu
suku-$\lambd$ khusus~$\num n$, yaitu \emph{numeral Church} untuk~$n$.
(Numeral Church dinamai menurut Alonzo Church.)

\begin{defn}
  Jika $n \in \Nat$, \emph{numeral Church} yang bersesuaian, $\num{n}$,
  merepresentasikan~$n$:
  \[
    \num{n} \ident \lambd[fx][f^n(x)]
  \]
  Di sini, $f^n(x)$ menyatakan hasil penerapan $f$ pada~$x$ sebanyak $n$
  kali. Sebagai contoh, $\num{0}$ adalah $\lambd[fx][x]$, sedangkan $\num{3}$
  adalah $\lambd[fx][f(f(f\,x))]$.
\end{defn}
  
Numeral Church $\num n$ dikodekan sebagai suku lambda yang merepresentasikan
suatu fungsi yang menerima dua argumen $f$ dan $x$, lalu mengembalikan
$f^n(x)$. Jelas bahwa numeral Church berbentuk normal.

Representasi bilangan asli dalam kalkulus lambda tentu hanya berguna jika kita
dapat melakukan komputasi dengannya. Melakukan komputasi dengan numeral Church
dalam kalkulus lambda berarti menerapkan suatu suku-$\lambd$~$F$ pada numeral
Church semacam itu, lalu mereduksi suku gabungan~$F\, \num n$ ke bentuk normal.
Jika suku tersebut selalu tereduksi ke bentuk normal dan bentuk normalnya selalu
berupa numeral Church~$\num m$, kita dapat menganggap keluaran komputasi itu
sebagai bilangan~$m$. Dengan demikian, kita dapat menganggap~$F$ mendefinisikan
suatu fungsi $f\colon \Nat \to \Nat$, yakni fungsi sedemikian sehingga
$f(n) = m$ jika dan hanya jika $F\, \num n \red \num m$. Berdasarkan sifat
Church--Rosser, bentuk normal, jika ada, bersifat tunggal. Jadi, jika
$F\, \num n \red \num m$, tidak mungkin terdapat suku lain berbentuk normal,
khususnya numeral Church lain, yang menjadi hasil reduksi $F \, \num n$.

Sebaliknya, jika diberikan suatu fungsi $f\colon \Nat \to \Nat$, kita dapat
bertanya apakah terdapat suku $F$ yang mendefinisikan~$f$ dengan cara ini.
Dalam hal itu kita mengatakan bahwa $F$ \emph{!!{lambda define}s}~$f$, dan
bahwa $f$ bersifat !!{lambda definable}. Gagasan ini dapat diperumum ke fungsi
berargumen banyak dan fungsi parsial.

\begin{defn}
Misalkan $f\colon \Nat^k \to \Nat$. Kita mengatakan bahwa suatu suku
lambda~$F$ \emph{!!{lambda define}s} $f$ jika, untuk semua
$n_0$, \dots,~$n_{k-1}$,
\[
F \, \num{n_0} \, \num{n_1} \dots \num{n_{k-1}} \red \num{f(n_0, n_1, \dots,
  n_{k-1})}
\]
apabila $f(n_0, \dots, n_{k-1})$ terdefinisi, sedangkan
$F \, \num{n_0} \, \num{n_1} \dots \num{n_{k-1}}$ tidak mempunyai bentuk
normal dalam kasus lainnya.
\end{defn}

Contoh yang sangat sederhana ialah fungsi-fungsi konstan. Suku $C_k \ident
\lambd[x][\num{k}]$ !!{lambda define}s fungsi $c_k\colon \Nat \to
\Nat$ sedemikian sehingga $c_k(n) = k$. Sebab,
$C_k \, \num n \ident (\lambd[x][\num{k}])\num n \redone \num{k}$ untuk
setiap~$n$.
Fungsi identitas !!{lambda defined} oleh $\lambd[x][x]$. Fungsi yang lebih
kompleks tentu lebih sukar didefinisikan dan sering kali menuntut banyak
kecerdikan. Karena itu, mungkin mengejutkan bahwa setiap fungsi yang dapat
dihitung bersifat !!{lambda definable}. Konversnya juga benar: jika suatu
fungsi bersifat !!{lambda definable}, fungsi itu dapat dihitung.

\end{document}
